%% AUTO-GENERATED by scripts/make_tables.py -- DO NOT EDIT BY HAND.
%% Source: results/research_questions.json
%% Status: STUB - no measurements available
\begin{table}[htbp]
\centering\small
\caption{프레임 예산의 단계별 분해를 평균이 아니라 $p_{95}$로 보고한다(식
\eqref{eq:frame-budget}). 총 $p_{95}$만으로는 재생이 끊긴다는 사실만 알 수 있고 어느
단계 때문인지는 알 수 없는데, 네 단계의 해결책은 서로 완전히 다르다. 평균 비용이 가장 큰
단계가 꼬리를 만드는 단계와 일치하지 않을 수 있으므로---예컨대 평소에는 저렴하지만 타일
예산을 간헐적으로 초과하는 rasteriser---귀속은 가장 느린 프레임들에 대해 계산한다.}
\label{tab:frame-budget-stages}
\begin{tabular}{@{}lrrrrrr@{}}
\toprule
설정 & 표면 & 정사영 & 방향 & rasteriser & 총 $p_{95}$ & 헤드룸 \\
 & (ms) & (ms) & (ms) & (ms) & (ms) & ($\times$) \\
\midrule
\NM & \NM & \NM & \NM & \NM & \NM & \NM \\
\bottomrule
\end{tabular}
\end{table}
